\begin{solucion}
 Sea $\alpha\in\mathbb{C}$ una raíz de $f(x)=x^{3}-3x+1$ entonces:
 \begin{align*}
    f(\alpha)&=\alpha^{3}-3\alpha+1=0 \quad \Rightarrow \quad \alpha^{3}=3\alpha-1 \qquad \text{(1.1)}\\
 \end{align*}
Queremos ver que $\beta=\alpha^{2}-2$ es raíz de $f$. Por lo que Sustituyendo en $f(x=\beta)$: 
\[
f(\beta)\;=\;(\alpha^{2}-2)^{3}-3(\alpha^{2}-2)+1.
\]
Expandiendo el término cúbico:
\[
(\alpha^{2}-2)^{3}\;=\;\alpha^{6}-6\alpha^{4}+12\alpha^{2}-8.
\]
Notemos que:
\begin{align*}
    \alpha^{4}&=\alpha\cdot\alpha^{3}\\
    &=\alpha(3\alpha-1) \qquad \text{usando (1.1)}\\
    &=3\alpha^{2}-\alpha,
\end{align*}
Por lo tanto:
\begin{align*}
    \alpha^{6}&=(\alpha^{3})^{2}\\
    &=(3\alpha-1)^{2}\\
    &=9\alpha^{2}-6\alpha+1\\
\end{align*}

Sustituyendo en la expresión de $f(\beta)$,
\begin{align*}
f(\beta)
  &=\bigl(\alpha^{6}-6\alpha^{4}+12\alpha^{2}-8\bigr)-3(\alpha^{2}-2)+1\\
  &=\bigl(9\alpha^{2}-6\alpha+1\bigr)-6(3\alpha^{2}-\alpha)
     +12\alpha^{2}-8-3\alpha^{2}+6+1\\
  &=9\alpha^{2}-6\alpha+1-18\alpha^{2}+6\alpha
     +12\alpha^{2}-8-3\alpha^{2}+7\\
  &=(9-18+12-3)\alpha^{2}+(-6\alpha+6\alpha)+(1-8+7)\\
  &=0+0+0 \;=\;0.
\end{align*}
De este modo, $f(\beta)=0$. Asi hemos demostrado que si $\alpha$ es raíz de $f(x)$, entonces
$\beta=\alpha^{2}-2$ también lo es. Por ultimo como $\beta$ es raíz entonces usando la misma relación concluimos que:
\[
\gamma=(\alpha^{2}-2)^{2}-2=(\beta)^{2}-2.
\]
también es raíz de $f(x)$.


Este resultado implica que las raíces de $f$ están relacionadas: a partir de una raíz $\alpha$,  
otra raíz se obtiene mediante la transformación $\alpha \mapsto \alpha^{2}-2$.  
Si denotamos $\gamma=\beta^{2}-2=(\alpha^{2}-2)^{2}-2$, esta sería la tercera raíz del polinomio.  
De hecho, aplicando una vez más la transformación, $\gamma^{2}-2$ vuelve a darnos $\alpha$,  
cerrando un \textbf{ciclo de longitud 3} entre las raíces.  
Esto concuerda con el hecho de que las tres raíces de $f(x)$ son \textbf{raíces conjugadas} entre sí,  
es decir, todas satisfacen el mismo polinomio mínimo irreducible sobre $\mathbb{Q}$  
[Clase 13, pág.\ 2].  
En otras palabras, $\alpha$, $\beta$ y $\gamma$ son conjugadas algebraicas.  
Concluimos que $\beta$ es raíz de $f$ siempre que $\alpha$ lo sea, como se quería demostrar.


% Este resultado implica que las raíces de $f$ están relacionadas:  
% a partir de una raíz $\alpha$, otra raíz se obtiene mediante la transformación
% $\alpha\mapsto\alpha^{2}-2$.  
% Si denotamos
% \[
% \gamma=\beta^{2}-2=(\alpha^{2}-2)^{2}-2,
% \]
% se obtiene la tercera raíz del polinomio.  
% Al aplicar de nuevo la transformación, $\gamma^{2}-2=\alpha$, cerrándose un ciclo de longitud
% $3$.  Esto concuerda con que las tres raíces de $f(x)$ son conjugadas entre sí, es decir,
% satisfacen el mismo polinomio mínimo irreducible sobre $\mathbb{Q}$ [Clase 13, pág.~2].  
% Concluimos que $\beta$ es raíz de $f$ siempre que $\alpha$ lo sea.

\end{solucion}